% !TeX TXS-program:compile = txs:///latex
% !TeX encoding = UTF-8
% !TeX spellcheck = <none>
\documentclass[uplatex,dvipdfmx,a4paper,10pt]{jsarticle}

\usepackage{amsmath,amssymb}
\usepackage{booktabs,array}
\usepackage[haranoaji,unicode]{pxchfon}
\usepackage[dvipdfmx]{graphicx,color}
\usepackage{url}

\topmargin=-2.0cm
\textheight=25.0cm
\textwidth=17.0cm
\oddsidemargin=-1.0cm
\evensidemargin=-1.0cm
\renewcommand{\baselinestretch}{0.95}
\hbadness=10000
\setlength{\columnsep}{0.75cm}
\makeatletter
\setlength{\@fptop}{0pt}
\setlength{\@dblfptop}{0pt}
\makeatother

\title{\textbf{進捗報告}}
\author{M1 TANG YU}
\date{2026年7月13日}

\begin{document}
\twocolumn[
\maketitle

\section*{調査目的と範囲}

SNNに導入する生理学的制約の根拠を明確にするため，既存の生理実験で報告されている網膜内機構の定量値を整理する。ヒトの視細胞・RGCトポグラフィ~\cite{curcio1990,curcioallen1990,watson2014}，マカクのH1応答および双極細胞経路~\cite{packer2002,smith2001,puthussery2013}，ならびにマーモセットRGCデータ~\cite{sridhar2025}を対象とし，各値の測定対象・条件・モデルへの直接写像の可否を区別してまとめた。

]

\clearpage
\onecolumn

\begingroup
\centering
\fontsize{7.3}{8.7}\selectfont
\setlength{\tabcolsep}{2.5pt}
\renewcommand{\arraystretch}{1.15}
\begin{tabular}{>{\raggedright\arraybackslash}p{3.10cm}
                >{\raggedright\arraybackslash}p{7.25cm}
                >{\centering\arraybackslash}p{1.35cm}
                >{\raggedright\arraybackslash}p{4.30cm}}
\toprule
機構／数値名 & 具体数値 & 所属物種 & 参考文献 \\
\midrule
視細胞・RGCの空間トポグラフィ
& coneの中心窩ピーク密度は平均199,000 cones/mm$^2$（100,000--324,000），総数は平均4.6百万（4.08--5.29百万）。RGC総数は0.7--1.5百万；midget RGC RF中心密度は偏心度依存であり，中心窩値はWatson式で33,163 deg$^{-2}$
& ヒト
& Curcio et al., 1990~\cite{curcio1990}；Curcio and Allen, 1990~\cite{curcioallen1990}；Watson, 2014~\cite{watson2014} \\

H1受容野の空間尺度・時間フィルタ
& combined RF径は4 mm偏心で122 $\mu$m，11 mm偏心で309 $\mu$m。時間応答の準線形fitでは，$\tau_1$は1,000 tdで3--5 ms，10 tdで30--50 ms，$\tau_2$は約4--5 ms；10$^\circ$刺激で共鳴は約40 Hz。いずれも刺激・順応依存のfit定数であり，コードの単一H1 $\tau$ ではない
& マカク
& Packer and Dacey, 2002~\cite{packer2002}；Smith et al., 2001~\cite{smith2001} \\

双極細胞のsustained／transient動態
& 集約BC stateに直接対応する共通のsustained／transient $\tau$ は未報告。元の形態学論文は時間応答を測定していない。macaque patch-clampはmagnocellular経路BCにNa$_\mathrm{V}$／Ca$_\mathrm{V}$／HCN電流が選択的に存在することを示すが，単一の光応答$\tau$は与えない
& マカク
& Puthussery et al., 2013~\cite{puthussery2013} \\

RGCの時間応答・spike精度
& M-cell dynamicsとspike精度の報告あり；単一時定数に相当する値は要確認
& マカク
& Benardete et al., 1999~\cite{benardete1999}；Uzzell et al., 2004~\cite{uzzell2004} \\

RGC参照データ（標本数）
& 摘出網膜2標本；naturalistic movie，白色雑音，反転gratingを収録
& マーモセット
& Sridhar et al., 2025~\cite{sridhar2025} \\

cone--H1局所接続数
& 各coneは網膜全域で約3--5個のH1，および3--5個のH2と接続
& マカク
& Wässle et al., 1989~\cite{wassle1989} \\

H1受容野径と偏心度
& combined RF diameterは4 mm偏心で122 $\mu$m，11 mm偏心で309 $\mu$m；約80\%は狭い中心と広く浅い成分で記述
& マカク
& Packer and Dacey, 2002~\cite{packer2002} \\

M/P中心尺度と周辺利得
& M center radiusは隣接Pの約2倍；M/Pともsurround/center integrated gain比は平均0.55；M contrast gainはPの約6倍
& マカク
& Croner and Kaplan, 1995~\cite{croner1995} \\

ON/OFFの空間・時間非対称
& ヒトdendritic diameterのON/OFF比はmidget $1.5\pm0.3$，parasol $1.3\pm0.1$；マカクparasol ON RFは約20\%大きく，TTP等は10--20\%短い
& ヒト／マカク
& Dacey and Petersen, 1992~\cite{dacey1992}；Chichilnisky and Kalmar, 2002~\cite{chichilnisky2002} \\

midget時間応答の偏心度依存
& STA TTPはfoveal $67\pm6$ ms，central $53\pm6$ ms，peripheral $37\pm4$ ms；biphasic indexは0.33，0.30，0.51付近
& マカク
& Sinha et al., 2017~\cite{sinha2017} \\

midget/parasolの局所収束
& foveal midgetは1 cone private line；典型的記録部位でmidget centerは約1--3 bipolar，parasol centerは少なくとも約30 bipolarを収束
& マカク
& Calkins et al., 1994~\cite{calkins1994}；Manookin et al., 2018~\cite{manookin2018} \\

ヒトmidget/parasol機能順序
& parasolはmidgetより時間RFが狭くbiphasic，空間RFが広く，response thresholdが低い；ONはOFFより広く低閾値
& ヒト
& Soto et al., 2020~\cite{soto2020} \\

parasol spike精度
& 強い高contrast刺激下でtrial間spike-time variabilityは最小約1 ms；spike countはsub-Poisson
& マカク
& Uzzell and Chichilnisky, 2004~\cite{uzzell2004} \\

自然視における微小眼球運動
& 動的自然場面でmicrosaccade rateは平均0.77 s$^{-1}$（0.5--1.1 s$^{-1}$）；典型的fixationは約400 ms
& ヒト
& Roberts et al., 2013~\cite{roberts2013} \\
\bottomrule
\end{tabular}
\par
\endgroup

\clearpage
\twocolumn

\begin{thebibliography}{99}
\small
\interlinepenalty=10000

\bibitem{curcio1990}
C. A. Curcio et al.,
``Human photoreceptor topography,''
\textit{Journal of Comparative Neurology}, 1990.
\url{https://doi.org/10.1002/cne.902920402}.

\bibitem{curcioallen1990}
C. A. Curcio and K. A. Allen,
``Topography of ganglion cells in human retina,''
\textit{Journal of Comparative Neurology}, 1990.
\url{https://doi.org/10.1002/cne.903000103}.

\bibitem{dacey1992}
D. M. Dacey and M. R. Petersen,
``Dendritic field size and morphology of midget and parasol ganglion cells of the human retina,''
\textit{Proceedings of the National Academy of Sciences}, 1992.
\url{https://doi.org/10.1073/pnas.89.20.9666}.

\bibitem{watson2014}
A. B. Watson,
``A formula for human retinal ganglion cell receptive field density as a function of visual field location,''
\textit{Journal of Vision}, 2014.
\url{https://doi.org/10.1167/14.7.15}.

\bibitem{verweij1996}
J. Verweij et al.,
``Horizontal cells feed back to cones by shifting the cone calcium-current activation range,''
\textit{Vision Research}, 1996.
\url{https://doi.org/10.1016/S0042-6989(96)00142-3}.

\bibitem{boycott1991}
B. B. Boycott and H. W\"assle,
``Cone bipolar cells and cone synapses in the primate retina,''
\textit{Visual Neuroscience}, 1991.
\url{https://doi.org/10.1017/S0952523800010932}.

\bibitem{hopkins1997}
J. M. Hopkins and B. B. Boycott,
``The cone synapses of cone bipolar cells of primate retina,''
\textit{Journal of Neurocytology}, 1997.
\url{https://doi.org/10.1023/A:1018504718282}.

\bibitem{benardete1999}
E. A. Benardete and E. Kaplan,
``The dynamics of primate M retinal ganglion cells,''
\textit{Visual Neuroscience}, 1999.
\url{https://doi.org/10.1017/S0952523899162151}.

\bibitem{uzzell2004}
V. J. Uzzell and E. J. Chichilnisky,
``Precision of spike trains in primate retinal ganglion cells,''
\textit{Journal of Neurophysiology}, 2004.
\url{https://doi.org/10.1152/jn.01171.2003}.

\bibitem{sridhar2025}
S. Sridhar and T. Gollisch,
``Dataset --- Marmoset retinal ganglion cell responses to naturalistic movies and spatiotemporal white noise,''
G-Node, 2025.
\url{https://doi.org/10.12751/g-node.t43ph1}.

\bibitem{wassle1989}
H. W\"assle, B. B. Boycott, and J. R\"ohrenbeck,
``Horizontal cells in the monkey retina: cone connections and dendritic network,''
\textit{European Journal of Neuroscience}, 1989.
\url{https://doi.org/10.1111/j.1460-9568.1989.tb00350.x}.

\bibitem{packer2002}
O. S. Packer and D. M. Dacey,
``Receptive field structure of H1 horizontal cells in macaque monkey retina,''
\textit{Journal of Vision}, 2002.
\url{https://doi.org/10.1167/2.4.1}.

\bibitem{smith2001}
V. C. Smith, J. Pokorny, B. B. Lee, and D. M. Dacey,
``Primate horizontal cell dynamics: an analysis of sensitivity regulation in the outer retina,''
\textit{Journal of Neurophysiology}, 85:545--558, 2001.
\url{https://doi.org/10.1152/jn.2001.85.2.545}.

\bibitem{puthussery2013}
T. Puthussery, S. Venkataramani, J. Gayet-Primo, R. G. Smith, and W. R. Taylor,
``Na$_\mathrm{V}$1.1 channels in axon initial segments of bipolar cells augment input to magnocellular visual pathways in the primate retina,''
\textit{Journal of Neuroscience}, 33:16045--16059, 2013.
\url{https://doi.org/10.1523/JNEUROSCI.1249-13.2013}.

\bibitem{croner1995}
L. J. Croner and E. Kaplan,
``Receptive fields of P and M ganglion cells across the primate retina,''
\textit{Vision Research}, 1995.
\url{https://doi.org/10.1016/0042-6989(94)E0066-T}.

\bibitem{chichilnisky2002}
E. J. Chichilnisky and R. S. Kalmar,
``Functional asymmetries in ON and OFF ganglion cells of primate retina,''
\textit{Journal of Neuroscience}, 2002.
\url{https://doi.org/10.1523/JNEUROSCI.22-07-02737.2002}.

\bibitem{sinha2017}
R. Sinha et al.,
``Cellular and circuit mechanisms shaping the perceptual properties of the primate fovea,''
\textit{Cell}, 2017.
\url{https://doi.org/10.1016/j.cell.2017.01.005}.

\bibitem{calkins1994}
D. J. Calkins et al.,
``M and L cones in macaque fovea connect to midget ganglion cells by different numbers of excitatory synapses,''
\textit{Nature}, 1994.
\url{https://doi.org/10.1038/371070a0}.

\bibitem{manookin2018}
M. B. Manookin, S. S. Patterson, and C. M. Linehan,
``Neural mechanisms mediating motion sensitivity in parasol ganglion cells of the primate retina,''
\textit{Neuron}, 2018.
\url{https://doi.org/10.1016/j.neuron.2018.02.006}.

\bibitem{soto2020}
F. Soto et al.,
``Efficient coding by midget and parasol ganglion cells in the human retina,''
\textit{Neuron}, 2020.
\url{https://doi.org/10.1016/j.neuron.2020.05.030}.

\bibitem{roberts2013}
J. A. Roberts, G. Wallis, and M. Breakspear,
``Fixational eye movements during viewing of dynamic natural scenes,''
\textit{Frontiers in Psychology}, 2013.
\url{https://doi.org/10.3389/fpsyg.2013.00797}.

\end{thebibliography}

\end{document}
